\title{xLSTM: Extended Long Short-Term Memory}

\begin{abstract}
The introduction of the constant error carousel and gating mechanisms in the 1990s established the foundational principles behind the Long Short-Term Memory (LSTM) architecture. LSTMs have since proven their longevity and played a significant role in pivotal advances in deep learning, most notably as the backbone of the first Large Language Models (LLMs). Nevertheless, the emergence of Transformer architectures, driven by the parallelism of self-attention mechanisms, ushered in a paradigm shift, surpassing LSTMs when applied at large scales. In this context, we revisit a fundamental inquiry: To what extent can language modeling be advanced by scaling up LSTMs to the order of billions of parameters, adopting advancements pioneered by current LLMs, while addressing established challenges of LSTMs? Our contributions begin with the proposal of exponential gating, incorporating suitable normalization and stabilization strategies. Additionally, we restructure the LSTM memory component, resulting in two variants: (i) sLSTM, which features a scalar-based memory and update scheme along with an updated approach to memory mixing, and (ii) mLSTM, designed with matrix memory enabling full parallelism and a covariance-based update methodology. These enhanced LSTM modules are integrated into residual block structures, forming xLSTM blocks, which are then aggregated to construct xLSTM architectures through residual stacking. The combination of exponential gating and revised memory designs endows xLSTM models with enhanced performance and scaling characteristics, demonstrating competitiveness with the latest Transformers and State Space Models in terms of both capability and scalability.\\
Code available at: \url{https://github.com/NX-AI/xlstm}
\end{abstract}

\section{Introduction}
\label{sec:intro}

The Long Short-Term Memory (LSTM) ideas~\citep{Hochreiter:91,Hochreiter:97nips,Hochreiter:97}, i.e., the constant error carousel and gating, were introduced to overcome the vanishing gradient problem of recurrent neural networks~\citep{Hochreiter:91,Hochreiter:00book}:

\begin{align}
\label{eq:lstm_idea}
  c_t \ = \ f_t \ c_{t-1} \ + \ 
          i_t \  z_t \ , \quad  
          h_t \ = \ o_t \ \psi({c_t}) \ . 
\end{align}

The constant error carousel operates by incrementally adjusting the cell state $c_{t-1}$ (highlighted in green) through cell inputs $z_t$, with this modification being regulated by the sigmoid gates (shown in blue). Specifically, the input gate $i_t$ and the forget gate $f_t$ manage this process of updating the cell state, whereas the output gate $o_t$ determines the release of information from the memory cell as the hidden state $h_t$. Afterward, the cell state undergoes normalization or compression via the function $\psi$, followed by gating at the output to produce the hidden state.

LSTMs have achieved widespread success across a multitude of fields \citep{Hochreiter:01,Hochreiter:07,Schmidhuber:15}, dominating the area of text generation prior to the emergence of Transformers in 2017~\citep{Vaswani:17}. Their capabilities have been thoroughly demonstrated on a variety of tasks involving sequential data, including but not limited to text generation~\citep{Graves:13,Karpathy:15blog}, handwriting synthesis~\citep{Graves:13}, sequence-to-sequence translation~\citep{Sutskever:14nips}, the evaluation of computer programs~\citep{Zaremba:14arxiv}, automated image captioning~\citep{Karpathy:15,Hossain:19}, source code generation~\citep{Karpathy:15blog}, rainfall-runoff prediction~\citep{Kratzert:18,Kratzert:19}, and hydrological forecasting for flood warnings~\citep{Nearing:24}. Within the domain of reinforcement learning, LSTMs continue to lead sequence modeling tasks, as seen in high-performing systems such as AlphaStar for StarCraft II~\citep{Vinyals:17short}, OpenAI Five for Dota 2~\citep{OpenAI:19}, and control models for magnetic confinement in nuclear fusion~\citep{Degrave:22short}. The strength of LSTMs lies in their ability to discover and represent abstract concepts, efficiently capturing semantic features and encoding them within their memory cells~\citep{Karpathy:15blog}, a property highlighted by findings such as number and syntax-specific neurons~\citep{Lakretz:19}, linguistic units~\citep{Bau:19}, and sentiment neurons~\citep{Radford:17arxiv}. LSTMs remain integral to important contemporary use cases~\citep{Degrave:22short,Nearing:24}, demonstrating their enduring value.

While LSTMs have achieved significant breakthroughs, they are constrained by three principal shortcomings: (i) Inability to revisit prior storage decisions. This is illustrated using the {\it Nearest Neighbor Search} scenario (see Appendix~\ref{sec:appNearestNeighborSearch}): Provided a reference vector, the model must process a sequence step-by-step to locate the vector most closely matching the reference and output its corresponding value at the sequence’s end. As depicted in the left panel of Figure~\ref{fig:lstmProblems}, the mean squared error highlights how LSTM models encounter difficulties updating previously stored values when encountering a more similar vector later; by contrast, our proposed xLSTM employs exponential gating to overcome this limitation. (ii) Constricted storage capacity, whereby information is limited to scalar representations within cell states. This restriction is demonstrated with {\it Rare Token Prediction}. The right panel in Figure~\ref{fig:lstmProblems} reports token prediction perplexity on Wikitext-103~\citep{Merity:17} for varying token frequencies. Due to their restricted storage abilities, LSTMs perform suboptimally when rare tokens are involved. The xLSTM approach circumvents this issue by leveraging matrix-based memory. (iii) Poor parallelizability as a result of memory interdependencies, specifically the recurrent links between hidden states across time steps, forcing strictly sequential computation.

These constraints in the LSTM framework have facilitated the transition to Transformers~\citep{Vaswani:17} within the field of language modeling. It raises the question: what advancements in language modeling accuracy and efficiency become feasible if one can circumvent these LSTM limitations and scale up LSTM architectures to match the parameter count of contemporary Large Language Models?

\section{Extended Long Short-Term Memory}
\label{sec:xLSTM}

In order to address the limitations inherent to LSTM, Extended Long Short-Term Memory (xLSTM) incorporates two primary enhancements to the foundational LSTM concept presented in Equation~\eqref{eq:lstm_idea}. Specifically, exponential gating and innovative memory mechanisms are introduced, thereby expanding the LSTM family with two distinct architectures: (i) sLSTM (refer to Section~\ref{sec:sLSTM}), characterized by scalar-based memory, scalar update dynamics, and memory intermixing, and (ii) mLSTM (see Section~\ref{sec:mLSTM}), distinguished by matrix memory structures and a covariance-based update rule utilizing outer products, which enables complete parallelization. Exponential gating serves as an improvement mechanism for both sLSTM and mLSTM. To maximize parallel computation, mLSTM foregoes traditional memory mixing, including recurrent interactions between hidden states. Both variants allow for the incorporation of multiple memory cells: in sLSTM, these cells can interact through memory mixing; moreover, sLSTM supports multiple heads with cell-level mixing restricted within each head, rather than across heads. This configuration, when coupled with exponential gating, introduces a novel approach to memory mixing. For mLSTM, utilizing multiple heads or cells produces equivalent operational outcomes.

When these LSTM extensions are embedded within residual block frameworks, they form xLSTM blocks, detailed in Section~\ref{sec:xLSTMblock}. Layering these xLSTM blocks in a residual manner creates xLSTM architectures, as explained in Section~\ref{sec:xLSTMarchitecure}. The structure and elements of the xLSTM architecture are depicted in Figure~\ref{fig:xlstm_sketch}.

\subsection{Review of the Long Short-Term Memory}
\label{sec:orgLSTM}

The foundational LSTM framework~\citep{Hochreiter:91,Hochreiter:97nips,Hochreiter:97} featured a scalar memory cell serving as both the computational and storage core. This design enabled effective mitigation of the vanishing gradient issue~\citep{Hochreiter:91,Hochreiter:00book}, primarily via the so-called constant error carousel mechanism that governs the cell state evolution. Within this architecture, the memory cell employs three distinct gating components: input, output, and forget gates—the latter of which was incorporated in subsequent developments by \citet{Gers:00}. The corresponding LSTM cell update equations at time step $t$ are specified as follows:

\begin{align}
c_t \ &= \  f_t \ c_{t-1} \ + \ i_t \ z_t &  & &\text{cell state} \\
h_t  \ &= \ o_t \ \tilde{h}_t \ , & \tilde{h}_t \ &= \ \psi \left( c_t\right)
  &\text{hidden state} \\
z_t \ &= \ \varphi \left( \tilde{z}_t \right) \ , 
  &\tilde{z}_t \ &=  \ \mathbf{w}^\top_{z} \ \mathbf{x}_t \ + \
  r_{z}  h_{t-1} \ + \  b_{z} \ \
  &\text{cell input} \\
i_t \ &= \ \sigma \left( \tilde{i}_t  \right) \ , 
  &\tilde{i}_t \ &= \ \mathbf{w}^\top_{i} \ \mathbf{x}_t \ + \
  r_{i}  \ h_{t-1} \ + \  b_{i} \ \
  &\text{input gate} \\
f_t \ &= \ \sigma \left(  \tilde{f}_t \right) \ , 
  &\tilde{f}_t \ &= \ \mathbf{w}^\top_{f} \ \mathbf{x}_t  \ + \
  r_{f}  \ h_{t-1} \ + \  b_{f} \ \
  &\text{forget gate} \\
o_t \ &= \ \sigma \left( \tilde{o}_t \right) \ , 
  &\tilde{o}_t  \ &= \ \mathbf{w}^\top_{o} \ \mathbf{x}_t \ + \
  r_{o}  \ h_{t-1} \ + \  b_{o} \ \
  &\text{output gate} 
\end{align}

The input weight vectors $\mathbf{w}_{z}$, $\mathbf{w}_{i}$, $\mathbf{w}_{f}$, and $\mathbf{w}_{o}$ denote the connections from the input $\mathbf{x}_t$ to the cell input, input gate, forget gate, and output gate, respectively. Similarly, $r_{z}$, $r_{i}$, $r_{f}$, and $r_{o}$ represent the recurrent weights from the hidden state $h_{t-1}$ to each corresponding component: cell input, input gate, forget gate, and output gate. The bias parameters, denoted $b_{z}$, $b_{i}$, $b_{f}$, and $b_{o}$, are associated with each of these gates and cell input.

The functions $\varphi$ and $\psi$ define activation for the cell input and the hidden state, often chosen as $\tanh$. Specifically, $\psi$ operates to bound the cell state, preventing it from diverging without constraint. All gating functions utilize the sigmoid activation, where $\sigma\left(x\right)=1/(1+\exp(-x))$.

Later developments involved aggregating several scalar memory cells into a vector $\mathbf{c}_t \in \mathbb{R}^{d}$, thereby enabling recurrent weight matrices $\mathbf{R} \in \mathbb{R}^{d \times d}$ that facilitate interaction among memory cell outputs~\citep{Greff:15}. Further explanation is available in Appendix~\ref{sec:apporgLSTM}. Investigations through ablation experiments established the necessity of each memory cell component for overall function~\citep{Greff:15}.

\subsection{sLSTM}
\label{sec:sLSTM}

To give LSTMs the capacity to adjust their storage choices, we propose exponential gates (red) supplemented by normalization and stabilization techniques. Specifically, the input and forget gates utilize exponential activation functions. For the purpose of normalization, we add a normalizer state, which accumulates the product of the input gate and the subsequent forget gates across future time steps.

The scalar sLSTM forward pass is:

\begin{align}
\label{eq:slstmforward}
c_t \ &= \  f_t \ c_{t-1} \ + \ i_t \ z_t & & 
 &\text{cell state} \\
n_t \ &= \  f_t \ n_{t-1} \ + \ i_t  & & 
 &\text{normalizer state} \\
h_t \ &= \ o_t \ \tilde{h}_t \ ,
  &\tilde{h}_t \ &= \ c_t / n_t
&\text{hidden state} \\
z_t \ &= \ \varphi \left( \tilde{z}_t \right) \ , 
  &\tilde{z}_t \ &=  \ \mathbf{w}^\top_{z} \ \mathbf{x}_t \ + \
  r_{z}  h_{t-1} \ + \  b_{z}
  &\text{cell input} \\
i_t \ &= \ \!\exp \left( \tilde{i}_t  \right) \ , 
  &\tilde{i}_t \ &= \ \mathbf{w}^\top_{i} \ \mathbf{x}_t \ + \
  r_{i}  \ h_{t-1} \ + \  b_{i}
  &\text{input gate} \\
f_t \ &= \ \sigma \left(  \tilde{f}_t \right) \ \text{OR} \
\!\exp \left(  \tilde{f}_t \right) \ ,
  &\tilde{f}_t \ &= \ \mathbf{w}^\top_{f} \ \mathbf{x}_t  \ + \
  r_{f}  \ h_{t-1} \ + \  b_{f}
  &\text{forget gate} \\
o_t \ &= \ \sigma \left( \tilde{o}_t \right) \ , 
  &\tilde{o}_t  \ &= \ \mathbf{w}^\top_{o} \ \mathbf{x}_t \ + \
  r_{o}  \ h_{t-1} \ + \  b_{o}
  &\text{output gate} 
\end{align}

We transfer the original LSTM gating techniques, i.e., input- and/or hidden-dependent gating plus bias term, to the new architectures. Exponential activation functions can lead to large values that cause overflows. Therefore, we stabilize gates with an additional state \mbox{$m_t$~\citep{Milakov:18arxiv}}:

\begin{align}
\label{eq:slstmstabil}
m_t \ &= \ \max \left( \log ( f_t ) + m_{t-1} , \log ( i_t ) \right) &\text{stabilizer state} \\
i'_t \ &= \ \!\exp \left( \log \left ( i_t \right) - m_t \right) \ = \ \!\exp \left( \tilde{i}_t - m_t \right) \ \, 
  &\text{stabil. input gate} \\
f'_t \ &= \ \!\exp \left( \log \left( f_t \right) + m_{t-1} - m_t \right) \ \, 
  &\text{stabil. forget gate}
\end{align}

Appendix~\ref{sec:appsLSTM} demonstrates that substituting $f_t$ with $f'_t$ and $i_t$ with $i'_t$ during the forward pass does not alter either the network's overall output or the gradients of the loss function relative to the parameters.

\paragraph{New Memory Mixing.} Similar to the classical LSTM, sLSTM supports multiple memory cells (refer to Appendix~\ref{sec:appsLSTM}). This allows for memory mixing facilitated by the recurrent connections $\mathbf{R}_{\mathbf{z}}$, $\mathbf{R}_{\mathbf{i}}$, $\mathbf{R}_{\mathbf{f}}$, and $\mathbf{R}_{\mathbf{o}}$, which propagate from the hidden state vector $\mathbf{h}$ to the memory cell input $\mathbf{z}$, as well as to the gate vectors $\mathbf{i}$, $\mathbf{f}$, and $\mathbf{o}$, respectively. A distinguishing feature in memory mixing arises from the use of exponential gating. In the revised sLSTM, multiple heads are possible, each enabling memory mixing internally but restricting interactions to within the head, not between different heads. By combining heads with exponential gating in sLSTM, a novel mechanism for mixing memory is introduced.

\subsection{mLSTM}
\label{sec:mLSTM}

To improve the memory capacity of LSTMs, the traditional scalar memory cell $c \in \mathbb{R}$ is replaced with a matrix $\mathbf{C} \in \mathbb{R}^{d \times d}$. Consequently, the retrieval process employs matrix multiplication. At each time step $t$, a vector pair is stored: a key $\mathbf{k}_t \in \mathbb{R}^d$ and a value $\mathbf{v}_t \in \mathbb{R}^d$ (adopting terminology from the Transformer architecture). At a later timestep $t+\tau$, the value $\mathbf{v}_t$ is recovered using a query vector $\mathbf{q}_{t+\tau} \in \mathbb{R}^d$. This configuration corresponds to the framework known as Bidirectional Associative Memories (BAMs) \citep{Kohonen:72,Anderson:72,Nakano:72,Anderson:77}.

The covariance update rule~\citep{Sejnowski:77,Dayan:91} for storing a key-value pair is

\begin{align}
\mathbf{C}_t \ &= \ \mathbf{C}_{t-1} \ + \ \mathbf{v}_{t} \ \mathbf{k}_t^\top \ .
\end{align}

We posit that a layer normalization is applied prior to the projection of inputs into keys and values, ensuring these representations maintain a zero mean. The covariance update mechanism achieves optimality~\citep{Dayan:91} with respect to maximizing the separability between recovered binary vectors, which, in turn, corresponds to achieving the highest possible signal-to-noise ratio. Enhanced separability can be attained if retrieval is restricted to pairwise relations, which results in quadratic computational complexity as observed in attention mechanisms~\citep{Krotov:16,Krotov:17,Ramsauer:21}. Notably, this covariance update is also functionally identical to that of Fast Weight Programmers~\citep{Schmidhuber:92ncfastweights,Schlag:21}, for which subsequent adaptations have introduced a fixed decay factor multiplied by $\mathbf{C}_{t-1}$, alongside a fixed learning rate that scales 
$\mathbf{v}_{t} \mathbf{k}_t ^\top$~\citep{Ba:16}. In alignment with this conceptual lineage, we embed the covariance update rule within the LSTM architecture—assigning the decay rate to the forget gate, the learning rate to the input gate, and utilizing the output gate to modulate the amplitude of the output vector.

For this matrix memory, the normalizer state is the weighted sum of key vectors, where each key vector is weighted by the input gate and all future forget gates. Again, the normalizer state keeps record of the strength of the gates. Since the dot product between query and normalizer state can be close to zero, we use the absolute value of this dot product and lower bound it by a threshold (typically 1.0) as done previously~\citep{Sun:23arxiv}. The mLSTM forward pass is:

\begin{align}
\label{eq:mlstm_recurrent_begin}
\mathbf{C}_t \ &= \  f_t \ \mathbf{C}_{t-1} \ + \ 
  i_t \ \mathbf{v}_t \ \mathbf{k}_t^\top & &  &\text{cell state} \\
\mathbf{n}_t \ &= \  f_t \ \mathbf{n}_{t-1} \ + \ 
  i_t \ \mathbf{k}_t & &  &\text{normalizer state} \\
\mathbf{h}_t  \ &= \ \mathbf{o}_t \ \odot \ \tilde{\mathbf{h}}_t
\ , \qquad \qquad 
\tilde{\mathbf{h}}_t \ = \   \mathbf{C}_t \mathbf{q}_t \ / \ 
  \max \left\{ |\mathbf{n}_t^\top \mathbf{q}_t|, 1 \right\} 
   &&&\text{hidden state} \\
\mathbf{q}_t \ &= \ \mathbf{W}_q \ \mathbf{x}_t \ + \ \mathbf{b}_q  & &  &\text{query input} \\
\mathbf{k}_t \ &= \ \frac{1}{\sqrt{d}} \mathbf{W}_k \ \mathbf{x}_t \ + \ \mathbf{b}_k  & &  &\text{key input} \\
\mathbf{v}_t \ &= \ \mathbf{W}_v \ \mathbf{x}_t \ + \ \mathbf{b}_v  & &  &\text{value input} \\
i_t \ &= \ \!\exp \!  \! \left( \tilde{i}_t  \right) \ , 
 \qquad \qquad \ \ \ \ \,
  \tilde{i}_t \ = \ \mathbf{w}^\top_{i} \ \mathbf{x}_t \ + \  b_{i} \ \ \ 
  &&&\text{input gate} \\
f_t \ &= \ \sigma \!  \left(  \tilde{f}_t \right) \ \text{OR} \ \!\exp \!  \! \left(  \tilde{f}_t \right) , \ \ \: \!
\tilde{f}_t \ = \ \mathbf{w}^\top_{f} \ \mathbf{x}_t  \ + \
  b_{f}
  &&& \text{forget gate} \\
\label{eq:mlstm_recurrent_end}
\mathbf{o}_t \ &= \ \sigma \left( \tilde{\mathbf{o}}_t \right) \ , \qquad \qquad \quad \ \ \ \,
  \tilde{\mathbf{o}}_t  \ = \ \mathbf{W}_{\mathbf{o}} \ \mathbf{x}_t \ + \
  \mathbf{b}_{\mathbf{o}}
  &&&\text{output gate} 
\end{align}

mLSTM retains the capacity to utilize several memory cells, similar to standard LSTM. In the context of mLSTM, employing multiple heads is functionally identical to using multiple memory cells, given the absence of memory mixing. To maintain stability in the exponential gating mechanisms of mLSTM, we adopt stabilization strategies akin to those implemented for sLSTM, as outlined in Equation~\ref{eq:slstmstabil}. Because mLSTM lacks memory mixing, its recurrence can be re-expressed in a parallel form. Additional information can be found in Appendix~\ref{sec:appmLSTM}.

\subsection{xLSTM Architecture}
\label{sec:xLSTMarch}

\paragraph{xLSTM Blocks.}
\label{sec:xLSTMblock}
The aim of an xLSTM block is to non-linearly encode historical input within a high-dimensional space, thereby enhancing the separation of distinct histories or contexts. This separation of sequence histories is essential for the accurate prediction of subsequent elements, such as the next token in a sequence. The design rationale is supported by Cover’s Theorem~\citep{Cover:65}, which argues that patterns embedded non-linearly in higher dimensional spaces are more likely to be separable by linear means compared to the original space. To this end, we examine two variants of residual block architectures: (i) The post up-projection residual block (as in Transformers), which first conducts a non-linear summary of the input within the initial feature space, followed by a linear transformation into an expanded high-dimensional space; this is then subjected to a non-linear activation and mapped back linearly to the initial space. The corresponding design is illustrated in the left panel of Figure~\ref{fig:Backbones}
and the third column of Figure~\ref{fig:xlstm_sketch}. 
A more comprehensive depiction is available in Figure~\ref{fig:appsLSTM_detailed}
in the appendix. (ii) The pre up-projection residual block (similar to State Space Models), which first employs a linear projection into a high-dimensional space,

compresses historical information in a non-linear fashion within a high-dimensional representation, followed by a linear transformation that restores it to the original space. When implementing an xLSTM block with an sLSTM, the up-projection component is positioned after the core unit. Conversely, in xLSTM blocks utilizing mLSTM, the up-projection is inserted before the core, as increasing dimensionality enhances the model's memory capacity. For more comprehensive illustrations, consult the left section of Figure~\ref{fig:Backbones}, the third column of Figure~\ref{fig:xlstm_sketch}, or refer to Figure~\ref{fig:appsLSTM_detailed} in the appendix.

\paragraph{xLSTM Architecture. }
\label{sec:xLSTMarchitecure}
The xLSTM architecture is composed by stacking modular units in a residual manner~\citep{Srivastava:15,He:16}. This design adopts the widely used pre-LayerNorm~\citep{Ba:16b} approach for residual connections, consistent with current Large Language Model practices. Refer to the final column of Figure~\ref{fig:xlstm_sketch} for visualization.

\subsection{Memory and Speed Considerations}

Unlike Transformers, xLSTM networks feature both linear computational complexity and constant memory requirements as sequence length increases. This compressive memory characteristic renders xLSTM particularly advantageous for industrial uses and deployment on edge devices.

In mLSTM, memory does not necessitate additional parameters but incurs substantial computational cost because of its reliance on a $d \times d$ matrix memory and $d \times d$ update operations. Here, memory capacity is balanced against computational complexity. However, since these matrix operations are highly amenable to GPU-based parallelization, their impact on wall clock time remains minimal.

Although mLSTM can be parallelized in a manner similar to methods such as FlashAttention~\citep{Dao:22,Dao:23} or GLA~\citep{Yang:23arxiv}, sLSTM presents a challenge for parallel execution due to the nature of its memory mixing through hidden-to-hidden connections. To address this, we engineered an efficient CUDA-based implementation with register-level GPU memory optimizations, achieving performance that is typically no more than twice as slow as mLSTM.

\section{Related Work}
\label{sec:related}

\paragraph{Linear Attention.} Numerous approaches have been proposed to reduce the quadratic computational demands of the Transformer's attention mechanism relative to context length, aiming for linear scalability instead. The Synthesizer circumvents direct token--token dependencies by generating synthetic attention weights~\citep{Tay:20arxiv}. The Linformer utilizes a low-rank matrix factorization to execute self-attention, achieving a linear approximation~\citep{Wang:20arxiv}. Linear Transformer restructures the attention process to attain linear complexity~\citep{Katharopoulos:20}. Performer employs positive orthogonal random features to derive a linear approximation of the attention softmax~\citep{Choromanski:21}. Additionally, methods such as Structured Global Convolution (SGConv)~\citep{Li:22} and Hyena Hierarchy~\citep{Poli:23} substitute conventional attention mechanisms with efficient long convolutional operations.

\paragraph{State Space Models.} In recent years, State Space Models (SSMs) have gained significant attention due to their linear complexity with respect to sequence length, making them strong contenders to Transformer-based architectures. Among the pioneering approaches was the Structured State Space sequence model (S4)~\citep{Gu:21}. This line of research has since expanded to include various extensions such as the Diagonal State Space (DSS) model~\citep{Gupta:22}, Gated State Space (GSS) models~\citep{Mehta:22}, S5 model~\citep{Smith:22}, Bidirectional Gated SSM (BiGS)~\citep{Wang:22}, H3 model~\citep{Fu:23}, as well as Mamba~\citep{Gu:24arxiv}.

\paragraph{Recurrent Neural Networks.} Recurrent Neural Networks (RNNs) have been proposed as alternatives to Transformers and attention mechanisms, owing to their ability to scale linearly with the sequence length. Deep Linear Recurrent Units (LRUs) have yielded encouraging outcomes in language modeling tasks~\citep{Orvieto:23, De:24arxiv}, alongside architectures such as Hierarchically Gated Linear RNNs (HGRN) \citep{Qin:23} and HGRN2~\citep{Qin:24arxiv}. Among notable RNN-based solutions for large language modeling is RWKV~\citep{Peng:23arxivshort,Peng:24arxivshort}, which demonstrates performance comparable to that of Transformer-based models.

\paragraph{Gating.}
Gating stands as a foundational concept within LSTM, and has subsequently been revisited and reformulated in a variety of more recent methodologies. Mechanisms based on gating principles appear across HGRN~\citep{Qin:23}, HGRN2~\citep{Qin:24arxiv}, Gated Linear Attention (GLA)~\citep{Yang:23arxiv}, Gated State Space (GSS) models~\citep{Mehta:22}, Bidirectional Gated SSM (BiGS)~\citep{Wang:22}, Moving Average Equipped Gated Attention (MEGA)~\citep{Ma:22}, RWKV~\citep{Peng:23arxivshort}, and Mamba~\citep{Gu:24arxiv}.

\paragraph{Covariance Update Rule.} To improve the capacity for memory storage, the mLSTM cell was augmented by introducing a matrix memory governed by a covariance update rule. Several other frameworks make use of related updating strategies, such as Fast Weight Programmers~\citep{Schmidhuber:92ncfastweights,Schlag:21}, RWKV-5 and RWKV-6~\citep{Peng:24arxivshort}, Retention~\citep{Sun:23arxiv}, Linear Transformer~\citep{Katharopoulos:20}, and HGRN2~\citep{Qin:24arxiv}.

\paragraph{Most Related.} The models that bear the closest resemblance to xLSTM in their conceptual design are Retention~\citep{Sun:23arxiv}, RWKV~\citep{Peng:23arxivshort,Peng:24arxivshort}, and HGRN2~\citep{Qin:24arxiv}. These systems similarly incorporate elements of matrix memory and/or gating mechanisms. Notably, unlike the new sLSTM architecture, these frameworks lack the capability for memory mixing. This property—memory mixing—provides an essential tool for addressing state tracking tasks. As a result, LSTM-based architectures surpass State Space Models (SSMs) and Transformers~\citep{Merrill:24,Deletang:23} in expressiveness. Efficient state tracking is a critical requirement for tasks such as code interpretation or following entities across extended narratives.

\paragraph{Residually Stacking Architectures.} The architecture of xLSTM, following the design of virtually all state-of-the-art deep learning models, is founded on the principle of residual stacking of modular blocks~\citep{Srivastava:15,He:16}.

Such a modular design paradigm is what allowed the development of highly deep convolutional networks~\citep{He:16} and Transformers~\citep{Vaswani:17}. Transformers in particular form the fundamental backbone for Large Language Models (LLMs), including but not limited to GPT-3~\citep{Brown:20short}, ChatGPT~\citep{Schulman:22short}, GPT-4~\citep{Achiam:23gpt4short}, Megatron-LM~\citep{Shoeybi:19}, Gopher~\citep{Rae:21short}, ERNIE 3.0 Titan~\citep{Wang:21arxivshort}, GLaM ~\citep{Du:21glamshort}, Chinese M6~\citep{Lin:21}, multilingual AlexaTM 20B~\citep{Soltan:22}, OPT~\citep{Zhang:22opt}, Chinchilla~\citep{Hoffmann:22short}, BLOOM~\citep{Scao:22arxivshort}, GLM-130B~\citep{Zeng:22arxivshort}, LaMDA~\citep{Thoppilan:22short}, PaLM~\citep{Chowdhery:22short}, Llama~\citep{Touvron:23arxiv}, and Gemini~\citep{GeminiTeam:23arxiv,Reid:24arxiv}.

\section{Experiments}
\label{sec:Exp}

This section presents an empirical assessment of xLSTM, benchmarking its performance against established approaches in the domain of language modeling. Section~\ref{sec:ExpSyntethic} explores xLSTM's properties through experiments on synthetic datasets. In Section~\ref{sec:ExpComparison}, we measure the validation perplexity achieved by a range of state-of-the-art language models, all trained using 15B tokens from the SlimPajama dataset~\citep{Soboleva:23}. Additionally, ablation analyses are carried out on xLSTM within this same dataset. We further examine how model performance evolves as a function of scaling, drawing parallels with the analyses by \citet{Kaplan:20} and \citet{Brown:20short}. 

Section~\ref{sec:ExpLanguage} encompasses a comprehensive language modeling investigation. Here, xLSTM is contrasted with top-performing models from Section~\ref{sec:ExpComparison}, each trained with an expanded dataset of 300B tokens from SlimPajama~\citep{Soboleva:23}. The evaluation proceeds in several stages: first, we probe model generalization to longer sequences; second, we report results on validation perplexity alongside downstream task accuracy~\citep{Sutawika:24}; third, we analyze performance across 571 distinct text domains featured in the PALOMA benchmark dataset~\citep{Magnusson:23arxivshort}; and finally, we revisit scaling analyses for all methods, this time using a dataset twenty times larger than previous experiments.

\label{sec:xlstm_blocks_notation}
Throughout all experiments, we denote the configuration xLSTM[$a$:$b$] to represent a composition of xLSTM blocks in the proportion $a/b$, with $a$ indicating mLSTM-based blocks and $b$ indicating sLSTM-based blocks. As an illustration, xLSTM[7:1] refers to an architecture where, out of a total of eight blocks, seven are mLSTM-based and one is sLSTM-based. When scaled to 48 blocks, this distribution results in 42 blocks utilizing mLSTM and 6 employing sLSTM. Additionally, in every experimental setup, mLSTM blocks are preceded by up-projection layers, while sLSTM blocks include post up-projection layers.

\subsection{Synthetic Tasks and Long Range Arena}
\label{sec:ExpSyntethic}
Initially, we examine how xLSTM’s novel exponential gating combined with memory mixing performs on tasks involving formal languages~\citep{Deletang:23}. Subsequently, we evaluate the capability of xLSTM’s advanced matrix memory structure in the Multi-Query Associative Recall task~\citep{Arora:23arxiv}. Lastly, the proficiency of xLSTM in managing extended input sequences is measured within the Long Range Arena benchmark~\citep{Tay:21}.

\paragraph{Evaluation of xLSTM's Exponential Gating and Memory Mixing Capabilities.} To assess the effectiveness of xLSTM's exponential gating combined with memory mixing—mechanisms designed to facilitate state tracking~\citep{Merrill:24, Merrill:23}—we adapt and expand the formal language experiments from \citet{Deletang:23}, incorporating support for multi-length training to enable length extrapolation. Appendix~\ref{sec:appExpSynthetic-formalLang} provides comprehensive task descriptions and additional empirical results. Our comparative analysis includes xLSTM alongside a range of models such as Transformers, State Space Models, and Recurrent Neural Networks. Accuracy is measured solely on those token predictions pertinent to the formal language tasks, normalized between 0 (random guessing) and 1 (complete correctness). Specifically, we examine 2-block variants of the following models: xLSTM[0:1] (only sLSTM), xLSTM[1:0] (only mLSTM), xLSTM[1:1], Llama, Mamba, RWKV, Retention, Hyena, LSTM, and a Transformer with LSTM blocks (LSTM (Block)). Figure~\ref{fig:formal-main} presents the experimental outcomes. Notably, architectures like Transformers and State Space Models that lack memory mixing—and, by extension, state tracking—are unable to solve tasks such as regular grammars (e.g., the parity task). This finding is consistent with prior work showing that, in terms of formal language recognition, Transformers and State Space Models possess fundamental limitations compared to RNNs~\citep{Merrill:24, Merrill:23, Deletang:23}.

\paragraph{Evaluation of xLSTM's Memory Capabilities Using Associative Recall Tasks.} This study investigates the matrix memory mechanism of xLSTM, focusing on its memory capacity within the Multi-Query Associative Recall task~\citep{Arora:23arxiv}. Each input sequence consists of key-value pairs drawn randomly from an extensive vocabulary, requiring the model to store and later retrieve these associations. To intensify the challenge compared to previous formulations, we scale up the number of key-value pairs to as many as 256 and expand the context length to 2048, thereby subjecting various models to more rigorous memory tests. Our comparison encompasses two-block configurations of Llama, Mamba, RWKV-5, RWKV-6, xLSTM[1:1], and xLSTM[1:0]. Performance is assessed through pair recall accuracy. Given that Transformers (such as Llama) possess memory capacities exponential relative to their coding dimension~\citep{Ramsauer:21}, they represent the benchmark for this task. The outcome, summarized in Figure~\ref{fig:mqar-main}, reveals that xLSTM[1:1] achieves the highest recall among non-Transformer architectures, excelling even when model size is limited. Notably, the inclusion of an sLSTM block does not detract from memory capability; in fact, it enhances performance, especially as the task reaches maximum difficulty with 256 pairs. Further findings, detailed in Appendix~\ref{sec:appExpSynthetic-mqar}, show through extrapolation studies that xLSTM's superior memory allows generalization to sequence contexts that exceed those encountered during training.

\paragraph{Test of xLSTM's Long Context Capabilities on Long Range Arena.} In order to evaluate how xLSTM manages extended sequences and substantial contextual information, we conduct a comparison of various techniques using the Long Range Arena~\citep{Tay:21}. Across the spectrum of tasks, xLSTM repeatedly exhibits robust results, implying that its architectural design is particularly adept at tackling diverse challenges associated with processing lengthy contexts.

For more details, see Appendix~\ref{sec:appExpSynthetic-lra}.

\subsection{Method Comparison and Ablation Study}
\label{sec:ExpComparison}

In order to investigate the central problem presented in this work—specifically, the scalability and performance of our newly introduced LSTM variants in language modeling tasks—we conduct training experiments using xLSTMs, Transformers, State Space Models, and additional approaches. All models are trained with 15B tokens derived from SlimPajama under an identical auto-regressive framework. Subsequently, we evaluate and contrast their performance on a shared validation dataset and carry out ablation analyses focused on the xLSTMs.

\paragraph{Comparison Between xLSTM and Alternative Approaches.} To facilitate a comparison, all models were trained using a dataset comprising 15B tokens sourced from SlimPajama~\citep{Soboleva:23}. Model performance was measured based on validation set perplexity. The approaches evaluated include: xLSTM (the proposed method), GPT-3 (Transformer architecture)~\citep{Brown:20short}, Llama (Transformer)~\citep{Touvron:23arxiv}, H3 (a State Space Model, SSM)~\citep{Fu:23}, Mamba (SSM)~\citep{Gu:24arxiv}, RWKV-4 (a Recurrent Neural Network, RNN)~\citep{Peng:23arxivshort}, RWKV-5 (RNN)~\citep{Peng:24arxivshort}, RWKV-6 (RNN)~\citep{Peng:24arxivshort}, GLA (a linear Transformer)~\citep{Yang:23arxiv}, HGRN (RNN)~\citep{Qin:23}, HGRN2 (RNN)~\citep{Qin:24arxiv}, RetNet (linear Transformer)~\citep{Sun:23arxiv}, Hyena (linear Transformer)~\citep{Poli:23}, xLSTM[1:0], and xLSTM[7:1] (refer to Section~\ref{sec:xlstm_blocks_notation} for definitions). All training was conducted using mixed precision; however, for models RWKV-5, RWKV-6, GLA, and HGRN2, mixed-precision training was applied without relying on PyTorch’s automated mixed precision functionality (see also Appendix Section~\ref{sec:appGenTrainProc}).

Methods are organized into three groups: (a) Transformers, (b) State Space Models (SSMs), and (c) Recurrent Neural Networks (RNNs) in conjunction with linear Transformers, the latter being linear approximations to the standard Transformer attention. All models were designed to match a GPT-3 model parameter count of 350M—specifically, using an embedding dimension of 1024 and 24 residual layers. Notably, only the GPT-3 model shares weights between the token and output embeddings, which results in a lower parameter count for this model compared to the others. According to the results presented in Table~\ref{tab:spaj15b_model_results}, xLSTM achieves the lowest validation perplexity among all compared models. Additional experimental details can be found in Appendix~\ref{sec:appExpComparison}. Furthermore, Figure~\ref{fig:temp_sclaw15B} illustrates the scaling properties observed in this setting, demonstrating that xLSTM continues to offer advantageous performance as model size increases.

\paragraph{Ablation Studies.}
Table~\ref{tab:spaj15b_model_results} and Figure~\ref{fig:temp_sclaw15B} present evidence that xLSTM delivers robust performance in language modeling tasks when trained with 15B tokens sourced from SlimPajama. To systematically investigate the transition from LSTM to xLSTM, we incrementally modify a standard LSTM architecture, introducing each key component of xLSTM in successive steps. Our process begins by embedding LSTM layers within pre-LayerNorm residual structures. Next, we add a post up-projection module. The final modification introduces exponential gating alongside matrix memory.

The results are shown in Table~\ref{tab:ablstudies} (top). 

Ablation experiments indicate that the notable gains in performance stem from the integration of exponential gating and matrix memory. Given the critical role of gating in RNNs and State Space Models, we further explore various gating configurations. As illustrated in Table~\ref{tab:ablstudies} (bottom), the results demonstrate that assigning learnability to each gate and conditioning them on the input yields a progressive improvement. Further analysis of the specific backbone modules is provided in Appendix~\ref{sec:appAblDetails}.

\subsection{xLSTM as Large Language Model}
\label{sec:ExpLanguage}

We conclude our investigation with comprehensive language modeling experiments, assessing xLSTM’s viability as a large language model (LLM). For this purpose, we scale up the training data, utilizing 300B tokens from SlimPajama—matching the token count adopted in prior works like Mamba~\citep{Gu:24arxiv} and Griffin~\citep{De:24arxiv}.

xLSTM's performance is benchmarked against RWKV-4, Llama, and Mamba—each serving as the representative for their respective architectural categories as discussed in Section~\ref{sec:ExpComparison}. RWKV-4 is chosen as the recurrent neural network (RNN) exemplar due to stable training precision settings not being established for RWKV-5, RWKV-6, and HGRN2 until after the commencement of our 300B-token training runs (see Appendix~\ref{sec:appGenTrainProc}).

Models are trained at various scales (125M, 350M, 760M, 1.3B) to examine their length extrapolation abilities, and their validation set performance is rigorously measured. Additional evaluations include their efficacy on downstream tasks, their language modeling performance across 571 text domains as delineated by the PALOMA benchmark, and an analysis of how each model adheres to scaling laws.

\paragraph{Sequence Length Extrapolation.}
\label{sec:spaj300B_ctx_extrapolation}
To begin, we investigate sequence length extrapolation abilities for large-scale (1.3B parameter) variants of xLSTM, RWKV-4, Llama, and Mamba. These models undergo training with a fixed context window of 2048 tokens, but are subsequently evaluated across much longer contexts, extending to 16384 tokens. Results are illustrated in Figure~\ref{fig:spaj300B_seq_extrapolation}. Notably, the xLSTM architecture displays consistently lower perplexity than its counterparts as context lengths increase.

\paragraph{Validation Perplexity and Downstream Tasks.}
\label{sec:spaj_downstream_eval}
Next, across each examined parameter scale, we assess xLSTM, RWKV-4, Llama, and Mamba in terms of next-token prediction accuracy on the SlimPajama validation set as well as performance on a suite of downstream tasks oriented toward evaluating common sense reasoning.

The third column in Table~\ref{tab:lmeval} presents the perplexities measured on the validation set for each approach. Among all evaluated model scales, both xLSTM[1:0] and xLSTM[7:1] exhibit the lowest validation perplexities. The remaining columns in Table~\ref{tab:lmeval} show the outcomes achieved on downstream tasks. xLSTM consistently outperforms competing methods across nearly all tasks and sizes, with the exception of a few instances in the ARC task where Mamba yields superior results. Further information can be found in Appendix~\ref{sec:appExpLanguage}.

\paragraph{Performance on PALOMA Language Tasks.} In the third evaluation, we assess the next token prediction accuracy of xLSTM, RWKV-4, Llama, and Mamba across varying model sizes, applying them to PALOMA language tasks~\citep{Magnusson:23arxivshort}. Performance is quantified via perplexity scores for predicting subsequent tokens in 571 diverse text domains, such as content from nytimes.com and the r/depression subreddit.

As depicted in Table~\ref{tab:ppleval}, the perplexity results are categorized by language modeling (the first seven columns) and by more granular domain benchmarks (the last five columns). Notably, xLSTM[1:0] demonstrates superior results to xLSTM[7:1] in these language tasks.

Compared to other models, xLSTM[1:0] achieves lower perplexity than Mamba in 568 of 571 domains (99.5\%), outperforms Llama in 486 of 571 domains (85.1\%), and surpasses RWKV-4 in 570 of 571 domains (99.8\%). Further specifics are provided in Appendix~\ref{sec:appExpLanguage}.

\paragraph{Scaling Laws.} Next, we investigate the observed power-law scaling patterns, which facilitate predictions about performance trends as model sizes increase~\citep{Kaplan:20,Brown:20short}. In Figure~\ref{fig:temp_sclaw300B}, we illustrate these scaling patterns. Although all examined models display analogous scaling characteristics, their performance levels differ due to varying offsets. RWKV-4 demonstrates the lowest performance among the group, with Llama and Mamba performing progressively better. Notably, xLSTM surpasses Mamba, maintaining a performance gap comparable to that between Mamba and Llama. This scaling analysis suggests that as model sizes grow, xLSTM is likely to maintain its performance advantage relative to both Transformer-based and State-Space architectures.

\textbf{Generation Times and Maximal Throughput.} We report the text generation latency in Figure~\ref{fig:inference_time_generation} and provide results for maximal throughput in Figure~\ref{fig:inference_time_decoding_throughput} (left), focusing on various xLSTM variants at the 1.3B parameter scale. Our evaluation benchmarks include similarly sized implementations of Mamba, Llama, and RWKV models from HuggingFace, with the Llama model utilizing a static key-value cache. During our experimental phase, neither complete cache compilation for the Transformer Model nor successful compilation of the Mamba model via \texttt{torch.compile} were available. For all text generation tasks, models were assessed using batch size 1 with a pre-fill of 16 tokens; this configuration particularly benefits the Transformer model. As depicted in Figure~\ref{fig:inference_time_generation}, xLSTM and the other recurrent models (Mamba and RWKV-4) demonstrate linear scaling with respect to sequence length, in contrast to Llama, which follows a quadratic scaling regime. Throughput analysis for Llama includes variations over batch size and prefill values. The right panel of Figure~\ref{fig:inference_time_decoding_throughput} highlights that, owing to constant memory requirements, xLSTM supports substantially larger batch sizes than Llama and thus attains the highest levels of throughput among the models considered.

\section{Limitations}

(i) Unlike mLSTM, the sLSTM architecture's memory mixing mechanism restricts the opportunity for parallel computation, thereby preventing highly efficient parallel processing. Despite this constraint, we have implemented a specialized CUDA kernel for sLSTM that currently operates at speeds just under double that of our parallelized mLSTM kernel.

(ii) The existing CUDA kernels for mLSTM have not been thoroughly optimized, resulting in runtimes roughly four times slower than those achieved with FlashAttention or the scan procedure utilized in Mamba. With targeted optimizations similar to those used for FlashAttention, substantial speed improvements in mLSTM kernels are possible.

(iii) mLSTM's matrix memory brings significant computational demands, as processing requires operations on $d \times d$ matrices. However, because memory updates and retrievals are parameter-free and amenable to conventional matrix operation parallelism, the additional processing time stemming from this complexity remains relatively negligible.

(iv) The initialization of forget gates warrants particular attention to avoid potential issues.

(v) Matrix memory size in mLSTM does not scale with sequence length, yet increasing sequence length risks overwhelming memory when working with longer contexts. Up to sequence lengths of 16k, this has not been problematic, as described in Section~\ref{sec:spaj300B_ctx_extrapolation}.

(vi) The high computational requirements associated with training large language models prevented exhaustive tuning of the architecture and hyperparameters, particularly with larger xLSTM variants. Achieving optimal xLSTM performance will likely require a thorough and systematic optimization process.

\section{Conclusion}
Our investigation provides a partial response to our central query: To what extent can LSTM architectures achieve in language modeling as their parameter count reaches billions? Up to this point, the evidence shows: ``At a minimum, xLSTM rivals established approaches such as Transformers and State Space Models''. We introduced enhancements to the traditional LSTM, transforming it into xLSTM through exponential gating combined with memory mixing as well as a novel memory configuration. Empirically, xLSTM consistently demonstrates strong performance in language modeling, outperforming or matching state-of-the-art models like Transformers and State Space Models. Analysis of scaling laws suggests that xLSTM models, when scaled up, could compete with leading Large Language Models primarily based on Transformer architectures. Furthermore, xLSTM may substantially influence domains such as Reinforcement Learning, Time Series Prediction, and modeling of physical phenomena.

\end{document}